\subsubsection{Periodic Coupled Boundaries}
\label{SEC:BC:Numerical:Periodic}

%###############################################################################

\devNotes{%
%
  \\\indent
%
  Currently, the translational vector is specified in grid scaled units.
  %
  This needs to be changed to grid units.
%
  }

%###############################################################################

Periodic coupled boundaries permit the flow exiting one boundary of the
computational domain to enter through a different boundary, and vice versa.
%
These two boundaries form a coupled pair of periodic boundary conditions.
%
Periodic boundaries fall into two categories of periodicity: translational and
rotational. A pair of periodic boundaries can be either translational,
rotational, or a combination of both.
%
Translational periodicity is quite useful for problems where the flow in one or
more directions can be treated as infinite, \forexample channel flow; details
and examples are discussed in
\sref{SEC:BC:Numerical:Periodic:Inputs:Translated}.
%
Rotational periodicity is useful for problems containing rotational symmetries,
\forexample a single blade from a turbomachinery rotor; details and examples are
discussed in \sref{SEC:BC:Numerical:Periodic:Inputs:Rotated}.

%%%%%%%%%%%%%%%%%%%%%%%%%%%%%%%%%%%%%%%%%%%%%%%%%%%%%%%%%%%%%%%%%%%%%%%%%%%%%%%%

We first start with some preliminary definitions and notation that will be used
for periodic coupled boundary conditions.
%
For all examples in this section, we will be working with a pair of periodic
boundary conditions that are specified in the input file using \BCinptext[i]{}
and \BCinptext[j]{}.
%
For the remainder of this section, every instance of \BCinptext[i]{} will have
the following two specified values
%
\begin{InputExample}
  \item[{\BCinptext[i]{Name = "Face Side A"}}]
  \item[{\BCinptext[i]{Set\_BC\_Value = 81}}]
\end{InputExample}
%
and similarly, \BCinptext[j] will have the following two specified values
%
\begin{InputExample}
  \item[{\BCinptext[j]{Name = "Periodic-2"}}]
  \item[{\BCinptext[j]{Set\_BC\_Value = 81}}]
\end{InputExample}
%
Note, the strings used here for the \inptext{Name} components of each boundary
condition have just been randomly created to illustrate that there is no
restriction that they be coordinated in any way.
%
In practice, it is recommended that they are coordinated just to make things
easier to understand, say something like \inptext{"Periodic-1A"} and
\inptext{"Periodic-1B"} if there might be more than one pair of periodic
boundary conditions.

%###############################################################################

\subsubsubsection{Things to Consider when Using Periodic Boundary Conditions}

\begin{itemize}
  %
  \item
    After all translation and rotation operations are performed on one side of a
    pair of coupled periodic boundaries, the linear grid components (linear cell
    vertices and faces) of the two boundaries need to result in a 1-to-1 match.
    %
    The algorithm for finding all of the matching pairs of vertices between two
    periodic boundaries works by computing the distance from the current vertex
    from the left boundary to all the vertices on the right boundary that have
    not yet been matched to a left vertex.
    %
    The matching right vertex is simply determined to be the one with the
    minimum distance to the current left vertex.
    %
    If the minimum distance is greater than 0.001
    \key[key:GridScaledUnits]{grid scaled units}[grid scaled units]
    (\thatis meters), an error will occur.

    For a pair of matching linear cell faces on periodic boundaries, the
    geometries for both faces of the pair need to match (\thatis both
    triangles, or both quadrilaterals), and the set of grid vertices for both
    faces need to create a matching set.
    %
    For example, let us say we have face $a$ defined by vertices $[1,2,3]$, and
    it was found to be a periodic match for face $b$ defined by vertices
    $[143,22,200]$.
    %
    For this to be a valid match, vertex 1 needs to match to one of the three
    vertices for face $b$; for this example we will say it matches to vertex
    200.
    %
    Vertex 2 then needs to be a match to either vertex 143 or 22; we will say 22
    for this example.
    %
    That means vertex 3 \emph{must} be a match to vertex 143, or else this is
    not a valid pair of matching faces.
    %
    Finally, the order of the matched vertices is not important here since after
    the faces are confirmed to be a valid match, the ordering will be used to
    find the relative rotation between the two faces.
  %
  \item
    Provided the previous bullet is satisfied such that the linear grid
    components of a pair of periodic boundaries are found to be an exact
    1-to-1 match, there are no other requirements on the topology, shape,
    or orientation of the boundaries.
    %
    Periodic boundaries can be high-order, curved surfaces, and they are not
    required to be planar or aligned with any coordinate direction\slash axis.
    %
    For example, a two-dimensional infinite turbine blade cascade could be
    simulated using only a single blade with translationally periodic boundaries
    along the curved geometric mid-passage lines above\slash below the blade.
  %
  \item
    For LES, make sure the distance between a pair of coupled
    periodic boundaries are separated by a distance that is sufficient enough to
    allow for turbulent structures to breakdown\slash dissipate\slash cascade
    into finer turbulent length scales along the direction of the periodicity.
    %
    If a pair of coupled periodic boundaries are too close to each other, the
    turbulent length scale will get artificially constrained in the periodic
    direction.
    %
    This will cause energy to get trapped in the larger\slash macro vortical
    structures, ultimately preventing a truly three-dimensional turbulent flow.
    
   %\FIXME{\\
   %This discusses breakdown, but don't periodic boundaries inhibit
   %\emph{growth} of turbulent structures?\\
   %Recommendation for spanwise width for LES?
   %Perhaps some multiple of the expected turbulent length scale?
   %}
  %
  \item
    Similar to the previous bullet, be careful with over-prescribing
    a simulation with too many periodic boundaries in such a way that the
    periodic boundary conditions change the flow dynamics.
    %
    The use of periodic boundary conditions is actually creating an infinite
    array of identical coupled simulations that are interacting with each other.
    %
    If the effect of this coupling is not well thought out, it can have
    unforeseen consequences ranging from unstable simulations that can be hard
    to diagnose the source of the instabilities, to questions on the validity
    \andor accuracy of a simulation and any conclusions that may have resulted.
     
    As mentioned in the previous bullet, having periodic boundaries too close to
    each other is an example of this, but another example is the fully periodic
    isentropic vortex problem\cite{Spiegel_2015}.
    %
    For this problem, 
    1) the Euler equations provide no physical dissipation;
    2) the high-order FR scheme provides minimal numerical dissipation; and
    3) the fully periodic boundaries prevent both numerical and machine
       roundoff errors from escaping the system.
    %
    With no means to dissipate the error or let it leave the system, all of the
    numerical and roundoff errors created by the simulation become trapped, and
    eventually these errors grow large enough to negatively effect the
    simulations with unstable\slash diverging solutions.
    %
    Replacing the periodic boundaries with characteristic boundary conditions
    completely eliminates all instabilities allowing us to demonstrate that
    the FR method achieves super accuracy for this problem.
    %
    Further details about this can be found in \linecite{Spiegel_2015}.
  %
\end{itemize}

%###############################################################################

\subsubsubsection{Specifying Coupled Periodic Boundary Conditions}
\label{SEC:BC:Numerical:Periodic:Inputs}

%###############################################################################

When specifying periodic boundary conditions, some input parameters are
commonly used with all such boundaries, and others are specific to
translational or rotational periodicity.
%
We discuss each in the sections that follow.

%###############################################################################

\FourSubsection{Inputs Common to All Periodic Boundaries}
\label{SEC:BC:Numerical:Periodic:Inputs:Common}

%This section discusses the inputs for specifying a pair of coupled periodic
%boundaries that are common to all periodic boundaries whether the periodicity
%is translational or rotational.
%
This section discusses inputs that can be used with both translational or
rotational periodic boundaries.
%
The \inptext{Periodic} subcomponent of \BCinpBC[nmlBC::BCInput]{} is used
to provide the information that will allow GFR to group two boundary conditions
into a matching pair of periodic boundaries, which in this example will be
\BCinptext[i]{} and \BCinptext[j]{}.
%
Both sides of a periodic pair must provide their respective information in the
input file.
%
%As you will see, the information required by opposing sides is the reverse\slash
%negative of each other, and could therefore probably be streamlined to reduce
%redundancies.
%%
%However, there was intentionality in how this was set up to help aid in error
%checking and making sure the user understands what they are doing in setting up
%periodic boundary conditions.
%
As you will see, the information required by opposing sides is the reverse\slash
negative of each other, and is used to aid error checking of the periodic
boundary setup.

%%%%%%%%%%%%%%%%%%%%%%%%%%%%%%%%%%%%%%%%%%%%%%%%%%%%%%%%%%%%%%%%%%%%%%%%%%%%%%%%

The two inputs used to define a matching periodic pair are
%
\begin{InputExample}
  \item[{\BCinptext[i]{Periodic\%Matching\_BC\_Name}}]
  \item[{\BCinptext[i]{Periodic\%Matching\_BC\_Idx}}]
\end{InputExample}
%
The values for these two inputs need to provide the respective information from
the opposing matching periodic boundary condition.
%
Using boundary conditions $i$ and $j$ defined at the beginning of
\sref{SEC:BC:Numerical:Periodic}, the inputs for boundary condition $i$ would
look like
%
\begin{InputExample}
  \item[{\BCinptext[i]{Name = "Face Side A"}}]
  \item[{\BCinptext[i]{Set\_BC\_Value = 81}}]
  \item[{\BCinptext[i]{Periodic\%Matching\_BC\_Name = "Periodic-2"}}]
  \item[{\BCinptext[i]{Periodic\%Matching\_BC\_Idx = }j}]
\end{InputExample}
%
and the inputs for boundary condition $j$ would look like
%
\begin{InputExample}
  \item[{\BCinptext[j]{Name = "Periodic-2"}}]
  \item[{\BCinptext[j]{Set\_BC\_Value = 81}}]
  \item[{\BCinptext[j]{Periodic\%Matching\_BC\_Name = "Face Side A"}}]
  \item[{\BCinptext[j]{Periodic\%Matching\_BC\_Idx = }i}]
\end{InputExample}
%
The remaining parts of \BCinptext{Periodic} are used to define the physical
orientation between the pair of periodic boundaries, and will be covered in the
next two sections.

%###############################################################################



%###############################################################################

\FourSubsection{Translated Periodic Boundary Condition}
\label{SEC:BC:Numerical:Periodic:Inputs:Translated}

Translated periodic boundary conditions are useful for problems where the flow
in one or more directions can be treated as infinite.
%
Examples primarily consist of two-dimensional geometries
(\forexample an airfoil)
that are converted to a three-dimensional body with an infinite span in the
third dimension.
%
Another example would be a 3D channel flow with translated periodic boundary
conditions in both the span-wise \emph{and} streamwise directions.
%
The periodicity in the streamwise direction allows for the flow to recirculate
indefinitely until it becomes a fully developed channel flow.
%
However, the streamwise periodicity requires a body force to drive the
streamwise bulk flow in the absence of a pressure differential generated by
traditional inflow\slash outflow boundary conditions.

%%%%%%%%%%%%%%%%%%%%%%%%%%%%%%%%%%%%%%%%%%%%%%%%%%%%%%%%%%%%%%%%%%%%%%%%%%%%%%%%

For each periodic boundary condition, the input vector
%
\begin{InputExample}
  \item[{\BCinpBC[nmlBC::BCInput/Periodic/Translation][i]%&
                 {Periodic\%Translation(1:3)}%
                 [Periodic\%Translation]}]
\end{InputExample}
%
provides the translation vector (in grid
units) that transforms the current boundary surface to align it exactly with its
corresponding periodic pair.
%
For a given periodic pair with translational periodicity, the following relation
must be satisfied:
%
\begin{InputExample}
  \item[{\BCinptext[i]{Periodic\%Translation(:) = -}
         \BCinptext[j]{Periodic\%Translation(:)}}]
\end{InputExample}

%%%%%%%%%%%%%%%%%%%%%%%%%%%%%%%%%%%%%%%%%%%%%%%%%%%%%%%%%%%%%%%%%%%%%%%%%%%%%%%%

For example, say you have an S-duct that is 5.0 inches long and drops 1.0 inch
vertically, and your grid file for this geometry is also provided in inches.
%
If you wanted to make this problem streamwise periodic, your inputs for boundary
condition $i$ would look like
%
\begin{InputExample}
  \item[{\BCinptext[i]{Periodic\%Translation(1:3) = [ 5.0, 1.0, 0.0]}}]
 %\item[{\BCinptext[i]{Periodic\%Translation(1) = 5.0}}]
 %\item[{\BCinptext[i]{Periodic\%Translation(2) = 1.0}}]
 %\item[{\BCinptext[i]{Periodic\%Translation(3) = 0.0}}]
\end{InputExample}
%
and the inputs for boundary condition $j$ would look like
%
\begin{InputExample}
  \item[{\BCinptext[j]{Periodic\%Translation(1:3) = [-5.0, -1.0, 0.0]}}]
 %\item[{\BCinptext[j]{Periodic\%Translation(1) = -5.0}}]
 %\item[{\BCinptext[j]{Periodic\%Translation(2) = -1.0}}]
 %\item[{\BCinptext[j]{Periodic\%Translation(3) = 0.0}}]
\end{InputExample}
%
Of course, an S-duct is generally not a problem that you would want to make
streamwise periodic.
%
It is merely used here to illustrate an example of the
periodic translational vector having components in multiple directions.

%###############################################################################

The elements on the boundary faces must be identical in order to have one-to-one
connectivity between cell faces.
%
When using translated periodic boundary conditions for large-eddy simulations,
one must make sure the extent of the periodic direction is long enough as to not
artificially constrain the turbulent length scale in the periodic direction.

%%%%%%%%%%%%%%%%%%%%%%%%%%%%%%%%%%%%%%%%%%%%%%%%%%%%%%%%%%%%%%%%%%%%%%%%%%%%%%%%



%###############################################################################

\FourSubsection{Rotated Periodic Boundary Condition}
\label{SEC:BC:Numerical:Periodic:Inputs:Rotated}

Rotated periodic boundary conditions are useful for problems containing
rotational symmetry.
%
This is particularly prevalent in turbomachinery flows where simulations of a
single periodic rotor blade can save significant computational expense compared
to simulations of the full 360-degree rotor wheel with dozens of rotor blades.
%
There are three subcomponent arrays of \BCinptext[i]{Periodic} that can be used
to specify the rotational transformation between a pair of rotated periodic
boundary conditions:
%
\begin{InputExample}
  \item[{\BCinpBC[nmlBC::BCInput/Periodic/RotationCenter][i]%&
                 {Periodic\%RotationCenter(1:3)}%
                 [Periodic\%RotationCenter]}]
  \item[{\BCinpBC[nmlBC::BCInput/Periodic/RotationAngle][i]%&
                 {Periodic\%RotationAngle(1:3)}%
                 [Periodic\%RotationAngle]}]
  \item[{\BCinpBC[nmlBC::BCInput/Periodic/RotationAxis][i]%&
                 {Periodic\%RotationAxis(1:4)}%
                 [Periodic\%RotationAxis]}]
\end{InputExample}
%
%One may use either \BCinptext[i]{Periodic\%RotationAngle} or
%\BCinptext[i]{Periodic\%RotationAxis} to define rotated periodic boundary
%conditions, but not both.
%
One may use either \inptext{RotationAngle} or \inptext{RotationAxis} to define
rotated periodic boundary conditions, but not both.
%
If both are specified, \inptext{RotationAxis} will be ignored and
\inptext{RotationAngle} will be used.
%
%For both methods, the input vector \BCinptext[i]{Periodic\%RotationCenter(1:3)}
%gives the rotation origin (in \key[key:GridUnits]{grid units}) used
%for defining the rotation angles between rotated periodic boundary conditions.
%
For both methods, the input vector \inptext{RotationCenter(1:3)} gives the
coordinates for the rotation origin (in \key[key:GridUnits]{grid units}) used
for defining the rotation angles between rotated periodic boundary conditions.

%%%%%%%%%%%%%%%%%%%%%%%%%%%%%%%%%%%%%%%%%%%%%%%%%%%%%%%%%%%%%%%%%%%%%%%%%%%%%%%%

The first way to specify the rotational transformation between a pair of rotated
periodic boundary conditions is through the input vector
\inptext{RotationAngle(1:3)}.
%\BCinptext[i]{Periodic\%RotationAngle(1:3)}.
%
This vector gives the rotation angles (in degrees) about each of the coordinate
axes from the current periodic boundary interface to the connecting periodic
boundary interface.
%
If rotating about more than one axis, the rotation is performed first about the
$x$-axis, then the $y$-axis, and finally the $z$-axis.
%
For a given periodic pair with rotational periodicity, the following relation
needs to be true:
%
\begin{InputExample}
  \item[\small{\BCinptext[i]{Periodic\%RotationAngle(:) = -}
               \BCinptext[j]{Periodic\%RotationAngle(:)}}]
\end{InputExample}

%%%%%%%%%%%%%%%%%%%%%%%%%%%%%%%%%%%%%%%%%%%%%%%%%%%%%%%%%%%%%%%%%%%%%%%%%%%%%%%%

The second way to specify the rotational transformation between a pair of
rotated periodic boundary conditions is through the input vector
\inptext{RotationAxis(1:4)}.
%\BCinptext[i]{Periodic\%RotationAxis(1:4)}.
%
The first 3 values of this array define a vector representing the axis of
rotation for this rotated periodic boundary condition, and the last value of the
array defines the rotation angle (in degrees) about this axis.
%
The magnitude of the vector defining the axis of rotation is not used; the sole
purpose of this vector is to define the direction of the rotation axis.
%
For a given periodic pair with rotational periodicity, one of the two following
sets of relations needs to be true:
%
\begin{InputExample}
  \item[\small{\BCinptext[i]{Periodic\%RotationAxis(1:3) = }
               \BCinptext[j]{Periodic\%RotationAxis(1:3)}}]
  \item[\small{\BCinptext[i]{Periodic\%RotationAxis(4) = -}
               \BCinptext[j]{Periodic\%RotationAxis(4)}}]
\end{InputExample}
%
or
%
\begin{InputExample}
  \item[\small{\BCinptext[i]{Periodic\%RotationAxis(1:3) = -}
               \BCinptext[j]{Periodic\%RotationAxis(1:3)}}]
  \item[\small{\BCinptext[i]{Periodic\%RotationAxis(4) = }
               \BCinptext[j]{Periodic\%RotationAxis(4)}}]
\end{InputExample}

%%%%%%%%%%%%%%%%%%%%%%%%%%%%%%%%%%%%%%%%%%%%%%%%%%%%%%%%%%%%%%%%%%%%%%%%%%%%%%%%

For example, say you have a rotor fan with 15 blades with the rotation axis
aligned with the Cartesian $x$-axis, and you would like to run a simulation of
just a single blade using rotational periodicity.
%
Using \inptext{RotationAngle}, your inputs for boundary condition $i$ would
look like
%
\begin{InputExample}
  \item[{\BCinptext[i]{Periodic\%RotationAngle(1:3) = [ 24.0, 0.0, 0.0]}}]
\end{InputExample}
%
and the inputs for boundary condition $j$ would look like
%
\begin{InputExample}
  \item[{\BCinptext[j]{Periodic\%RotationAngle(1:3) = [-24.0, 0.0, 0.0]}}]
\end{InputExample}
%
If you wanted to use \inptext{RotationAxis} instead, your inputs for boundary
condition $i$ would look like
%
\begin{InputExample}
  \item[{\BCinptext[i]{Periodic\%RotationAxis(1:4) = [ 1.0, 0.0, 0.0, 24.0]}}]
\end{InputExample}
%
and the inputs for boundary condition $j$ would look like
%
\begin{InputExample}
  \item[{\BCinptext[j]{Periodic\%RotationAxis(1:4) = [ 1.0, 0.0, 0.0,-24.0]}}]
\end{InputExample}



%###############################################################################

\subsubsubsection{Past Examples of Periodic Boundary Conditions}
\label{SEC:BC:Numerical:Periodic:Examples}

These are examples taken from GFR input files that have been used in the past to
run simulations for problems including periodic boundaries.
%
Note that the array indices used for \inptext{BC\_Input} are simply copied over
from the input files, meaning each index is relative to the other boundary
conditions that also existed in the input file it was taken from.
%
Part of the intent behind keeping these array indices as they originally were in
their respective input file is to demonstrate that the previously used variable
indices of $i$ and $j$ are can represent any array element of
\inptext{BC\_Input} provided that they link to each other when they form a pair
of periodic boundaries.
%
\begin{itemize}
  \item Streamwise and spanwise periodic channel flow
   %\begin{nmldescitem}[escapechar=@]
   %\begin{nmldescitem}[fontupper=\ttfamily\footnotesize,
    \begin{nmldescitem}[fontupper=\ttfamily\scriptsize,
                        listing options={escapeinside={!@}{!@}}]
&nml_bc
  BC_Input(1)%Name = "INFLOW"
  BC_Input(1)%Set_BC_Value = 81
  BC_Input(1)%Periodic%Matching_BC_Name = "OUTFLOW"
  BC_Input(1)%Periodic%Matching_BC_Idx = 2
  BC_Input(1)%Periodic%Translation(1) = 0.036576 !@\vspace{5pt}!@
  BC_Input(2)%Name = "OUTFLOW"
  BC_Input(2)%Set_BC_Value = 81
  BC_Input(2)%Periodic%Matching_BC_Name = "INFLOW"
  BC_Input(2)%Periodic%Matching_BC_Idx = 1
  BC_Input(2)%Periodic%Translation(1) = -0.036576 !@\vspace{5pt}!@
  BC_Input(5)%Name = "ZMIN"
  BC_Input(5)%Set_BC_Value = 81
  BC_Input(5)%Periodic%Matching_BC_Name = "ZMAX"
  BC_Input(5)%Periodic%Matching_BC_Idx = 6
  BC_Input(5)%Periodic%Translation(3) = 0.018288 !@\vspace{5pt}!@
  BC_Input(6)%Name = "ZMAX"
  BC_Input(6)%Set_BC_Value = 81
  BC_Input(6)%Periodic%Matching_BC_Name = "ZMIN"
  BC_Input(6)%Periodic%Matching_BC_Idx = 5
  BC_Input(6)%Periodic%Translation(3) = -0.018288
/
    \end{nmldescitem}
   %\FINALFIX{\newpage}
  \item Single rotor blade from spacecraft cabin ventilation
    fan\cite{Spiegel_2025}
    \begin{itemize}[before=\small]
      \item Rotor consists of nine blades spinning about the $x$-axis resulting
        in a rotationally periodic domain that is rotated 40\textdegree{} about
        the $x$-axis.
      \item This grid has two pairs of rotationally periodic boundaries: one for
        the primary computational domain, and a second for a sponge region
        appended onto the regular outflow boundary to damp out non-physical
        pressure waves from reflecting off the outflow boundary.
        The periodic boundary is split in this way for convenience, but could
        be combined.
    \end{itemize}
   %\begin{nmldescitem}[escapechar=@]
   %\begin{nmldescitem}[fontupper=\ttfamily\footnotesize,
    \begin{nmldescitem}[fontupper=\ttfamily\scriptsize,
                        listing options={escapeinside={!@}{!@}}]
&nml_bc
  BC_Input(8)%Name = "ThetaMin"
  BC_Input(8)%Set_BC_Value = 81
  BC_Input(8)%Periodic%Matching_BC_Name = "ThetaMax"
  BC_Input(8)%Periodic%Matching_BC_Idx = 9
  BC_Input(8)%Periodic%RotationAngle(1) = 40.0 !@\vspace{5pt}!@
  BC_Input(9)%Name = "ThetaMax"
  BC_Input(9)%Set_BC_Value = 81
  BC_Input(9)%Periodic%Matching_BC_Name = "ThetaMin"
  BC_Input(9)%Periodic%Matching_BC_Idx = 8
  BC_Input(9)%Periodic%RotationAngle(1) = -40.0 !@\vspace{5pt}!@
  BC_Input(10)%Name = "ThetaMinSponge"
  BC_Input(10)%Set_BC_Value = 81
  BC_Input(10)%Periodic%Matching_BC_Name = "ThetaMaxSponge"
  BC_Input(10)%Periodic%Matching_BC_Idx = 11
  BC_Input(10)%Periodic%RotationAngle(1) = 40.0 !@\vspace{5pt}!@
  BC_Input(11)%Name = "ThetaMaxSponge"
  BC_Input(11)%Set_BC_Value = 81
  BC_Input(11)%Periodic%Matching_BC_Name = "ThetaMinSponge"
  BC_Input(11)%Periodic%Matching_BC_Idx = 10
  BC_Input(11)%Periodic%RotationAngle(1) = -40.0
/
    \end{nmldescitem}
  \item Ceramic Matrix Composite 9\% thickness (CMC9) blade
    cascade\cite{Giel_2020} with spanwise and vertical
    periodicity\cite{Miki_2022}
   %\begin{nmldescitem}[escapechar=@]
   %\begin{nmldescitem}[fontupper=\ttfamily\footnotesize,
    \begin{nmldescitem}[fontupper=\ttfamily\scriptsize,
                        listing options={escapeinside={!@}{!@}}]
&nml_bc
  BC_Input(3)%Name = "periodic_yp13"
  BC_Input(3)%Set_BC_Value = 81
  BC_Input(3)%Periodic%Matching_BC_Name = "periodic_ym"
  BC_Input(3)%Periodic%Matching_BC_Idx = 4
  BC_Input(3)%Periodic%Translation(2) = -0.13 !@\vspace{5pt}!@
  BC_Input(4)%Name = "periodic_ym"
  BC_Input(4)%Set_BC_Value = 81
  BC_Input(4)%Periodic%Matching_BC_Name = "periodic_yp13"
  BC_Input(4)%Periodic%Matching_BC_Idx = 3
  BC_Input(4)%Periodic%Translation(2) = 0.13 !@\vspace{5pt}!@
  BC_Input(5)%Name = "periodic_zp013"
  BC_Input(5)%Set_BC_Value = 81
  BC_Input(5)%Periodic%Matching_BC_Name = "periodic_z0"
  BC_Input(5)%Periodic%Matching_BC_Idx = 6
  BC_Input(5)%Periodic%Translation(3) = -0.052 !@\vspace{5pt}!@
  BC_Input(6)%Name = "periodic_z0"
  BC_Input(6)%Set_BC_Value = 81
  BC_Input(6)%Periodic%Matching_BC_Name = "periodic_zp013"
  BC_Input(6)%Periodic%Matching_BC_Idx = 5
  BC_Input(6)%Periodic%Translation(3) = 0.052
/
    \end{nmldescitem}
  \item Taylor-Couette flow\cite{Wang_2021} with streamwise periodicity
   %\begin{nmldescitem}[escapechar=@]
   %\begin{nmldescitem}[fontupper=\ttfamily\footnotesize,
    \begin{nmldescitem}[fontupper=\ttfamily\scriptsize,
                        listing options={escapeinside={!@}{!@}}]
&nml_bc
  BC_Input(3)%Name = "Inflow"
  BC_Input(3)%Set_BC_Value = 81
  BC_Input(3)%Periodic%Matching_BC_Name = "Outflow"
  BC_Input(3)%Periodic%Matching_BC_Idx = 4
  BC_Input(3)%Periodic%Translation(1) = 6.283185307 !@\vspace{5pt}!@
  BC_Input(4)%Name = "Outflow"
  BC_Input(4)%Set_BC_Value = 81
  BC_Input(4)%Periodic%Matching_BC_Name = "Inflow"
  BC_Input(4)%Periodic%Matching_BC_Idx = 3
  BC_Input(4)%Periodic%Translation(1) = -6.283185307
/
    \end{nmldescitem}
\end{itemize}

%###############################################################################

%\begin{itemize}
%  \item Streamwise and spanwise periodic channel flow
%    \begin{InputExample}[before=\footnotesize]
%      %
%      \item[{\BCinptext[\texttt{1}]{Name = "INFLOW"}}]
%      \item[{\BCinptext[\texttt{1}]{Set\_BC\_Value = 81}}]
%      \item[{\BCinptext[\texttt{1}]{Periodic\%Matching\_BC\_Name = "OUTFLOW"}}]
%      \item[{\BCinptext[\texttt{1}]{Periodic\%Matching\_BC\_Idx = 2}}]
%      \item[{\BCinptext[\texttt{1}]{Periodic\%Translation(1) = 0.036576}}]
%      %
%      \item[{\BCinptext[\texttt{2}]{Name = "OUTFLOW"}}]\vspace{5pt}
%      \item[{\BCinptext[\texttt{2}]{Set\_BC\_Value = 81}}]
%      \item[{\BCinptext[\texttt{2}]{Periodic\%Matching\_BC\_Name = "INFLOW"}}]
%      \item[{\BCinptext[\texttt{2}]{Periodic\%Matching\_BC\_Idx = 1}}]
%      \item[{\BCinptext[\texttt{2}]{Periodic\%Translation(1) = -0.036576}}]
%      %
%      \item[{\BCinptext[\texttt{5}]{Name = "ZMIN"}}]\vspace{5pt}
%      \item[{\BCinptext[\texttt{5}]{Set\_BC\_Value = 81}}]
%      \item[{\BCinptext[\texttt{5}]{Periodic\%Matching\_BC\_Name = "ZMAX"}}]
%      \item[{\BCinptext[\texttt{5}]{Periodic\%Matching\_BC\_Idx = 6}}]
%      \item[{\BCinptext[\texttt{5}]{Periodic\%Translation(3) = 0.018288}}]
%      %
%      \item[{\BCinptext[\texttt{6}]{Name = "ZMAX"}}]\vspace{5pt}
%      \item[{\BCinptext[\texttt{6}]{Set\_BC\_Value = 81}}]
%      \item[{\BCinptext[\texttt{6}]{Periodic\%Matching\_BC\_Name = "ZMIN"}}]
%      \item[{\BCinptext[\texttt{6}]{Periodic\%Matching\_BC\_Idx = 5}}]
%      \item[{\BCinptext[\texttt{6}]{Periodic\%Translation(3) = -0.018288}}]
%      %
%    \end{InputExample}
%  \item Single rotor blade from spacecraft cabin ventilation
%    fan\cite{Spiegel_2025}
%    \begin{itemize}
%      \item Rotor consists of nine blades spinning about the $x$-axis resulting
%        in a rotationally periodic domain that is rotated 40\textdegree{} about
%        the $x$-axis.
%      \item This grid has two pairs of rotationally periodic boundaries: one for
%        the primary computational domain, and a second for a sponge region
%        appended onto the regular outflow boundary to damp out non-physical
%        pressure waves from reflecting off the outflow boundary.
%    \end{itemize}
%    \begin{InputExample}[before=\footnotesize]
%      %
%      \item[{\BCinptext[\texttt{8}]{Name = "ThetaMin"}}]
%      \item[{\BCinptext[\texttt{8}]{Set\_BC\_Value = 81}}]
%      \item[{\BCinptext[\texttt{8}]{Periodic\%Matching\_BC\_Name = "ThetaMax"}}]
%      \item[{\BCinptext[\texttt{8}]{Periodic\%Matching\_BC\_Idx = 9}}]
%      \item[{\BCinptext[\texttt{8}]{Periodic\%RotationAngle(1) = 40.0}}]
%      %
%      \item[{\BCinptext[\texttt{9}]{Name = "ThetaMax"}}]\vspace{5pt}
%      \item[{\BCinptext[\texttt{9}]{Set\_BC\_Value = 81}}]
%      \item[{\BCinptext[\texttt{9}]{Periodic\%Matching\_BC\_Name = "ThetaMin"}}]
%      \item[{\BCinptext[\texttt{9}]{Periodic\%Matching\_BC\_Idx = 8}}]
%      \item[{\BCinptext[\texttt{9}]{Periodic\%RotationAngle(1) = -40.0}}]
%      %
%      \item[{\BCinptext[\texttt{10}]{Name = "ThetaMinSponge"}}]\vspace{5pt}
%      \item[{\BCinptext[\texttt{10}]{Set\_BC\_Value = 81}}]
%      \item[{\BCinptext[\texttt{10}]{Periodic\%Matching\_BC\_Name = "ThetaMaxSponge"}}]
%      \item[{\BCinptext[\texttt{10}]{Periodic\%Matching\_BC\_Idx = 11}}]
%      \item[{\BCinptext[\texttt{10}]{Periodic\%RotationAngle(1) = 40.0}}]
%      %
%      \item[{\BCinptext[\texttt{11}]{Name = "ThetaMaxSponge"}}]\vspace{5pt}
%      \item[{\BCinptext[\texttt{11}]{Set\_BC\_Value = 81}}]
%      \item[{\BCinptext[\texttt{11}]{Periodic\%Matching\_BC\_Name = "ThetaMinSponge"}}]
%      \item[{\BCinptext[\texttt{11}]{Periodic\%Matching\_BC\_Idx = 10}}]
%      \item[{\BCinptext[\texttt{11}]{Periodic\%RotationAngle(1) = -40.0}}]
%      %
%    \end{InputExample}
%  \item Ceramic Matrix Composite 9\% thickness (CMC9) blade\cite{Giel_2020} with
%    spanwise and vertical periodicity\cite{Miki_2022}
%    \begin{InputExample}[before=\footnotesize]
%      %
%      \item[{\BCinptext[\texttt{3}]{Name = "periodic\_yp13"}}]
%      \item[{\BCinptext[\texttt{3}]{Set\_BC\_Value = 81}}]
%      \item[{\BCinptext[\texttt{3}]{Periodic\%Matching\_BC\_Name = "periodic\_ym"}}]
%      \item[{\BCinptext[\texttt{3}]{Periodic\%Matching\_BC\_Idx = 4}}]
%      \item[{\BCinptext[\texttt{3}]{Periodic\%Translation(2) = -0.13}}]
%      %
%      \item[{\BCinptext[\texttt{4}]{Name = "periodic\_ym"}}]\vspace{5pt}
%      \item[{\BCinptext[\texttt{4}]{Set\_BC\_Value = 81}}]
%      \item[{\BCinptext[\texttt{4}]{Periodic\%Matching\_BC\_Name = "periodic\_yp13"}}]
%      \item[{\BCinptext[\texttt{4}]{Periodic\%Matching\_BC\_Idx = 3}}]
%      \item[{\BCinptext[\texttt{4}]{Periodic\%Translation(2) = 0.13}}]
%      %
%      \item[{\BCinptext[\texttt{5}]{Name = "periodic\_zp013"}}]\vspace{5pt}
%      \item[{\BCinptext[\texttt{5}]{Set\_BC\_Value = 81}}]
%      \item[{\BCinptext[\texttt{5}]{Periodic\%Matching\_BC\_Name = "periodic\_z0"}}]
%      \item[{\BCinptext[\texttt{5}]{Periodic\%Matching\_BC\_Idx = 6}}]
%      \item[{\BCinptext[\texttt{5}]{Periodic\%Translation(3) = -0.052}}]
%      %
%      \item[{\BCinptext[\texttt{6}]{Name = "periodic\_z0"}}]\vspace{5pt}
%      \item[{\BCinptext[\texttt{6}]{Set\_BC\_Value = 81}}]
%      \item[{\BCinptext[\texttt{6}]{Periodic\%Matching\_BC\_Name = "periodic\_zp013"}}]
%      \item[{\BCinptext[\texttt{6}]{Periodic\%Matching\_BC\_Idx = 5}}]
%      \item[{\BCinptext[\texttt{6}]{Periodic\%Translation(3) = 0.052}}]
%      %
%    \end{InputExample}
%  \item Taylor-Couette flow\cite{Wang_2021} with streamwise periodicity
%    \begin{InputExample}[before=\footnotesize]
%      %
%      \item[{\BCinptext[\texttt{3}]{Name = "Inflow"}}]
%      \item[{\BCinptext[\texttt{3}]{Set\_BC\_Value = 81}}]
%      \item[{\BCinptext[\texttt{3}]{Periodic\%Matching\_BC\_Name = "Outflow"}}]
%      \item[{\BCinptext[\texttt{3}]{Periodic\%Matching\_BC\_Idx = 4}}]
%      \item[{\BCinptext[\texttt{3}]{Periodic\%Translation(1) = 6.283185307}}]
%      %
%      \item[{\BCinptext[\texttt{4}]{Name = "Outflow"}}]\vspace{5pt}
%      \item[{\BCinptext[\texttt{4}]{Set\_BC\_Value = 81}}]
%      \item[{\BCinptext[\texttt{4}]{Periodic\%Matching\_BC\_Name = "Inflow"}}]
%      \item[{\BCinptext[\texttt{4}]{Periodic\%Matching\_BC\_Idx = 3}}]
%      \item[{\BCinptext[\texttt{4}]{Periodic\%Translation(1) = -6.283185307}}]
%      %
%    \end{InputExample}
%\end{itemize}

%###############################################################################

\IDEA{Each example should include a schematic of the labeled boundary
surfaces to help visualize the periodicity.}

%###############################################################################



%###############################################################################
%###############################################################################
%###############################################################################
%###############################################################################
%###############################################################################

